% ECONOMICS_PROBLEM: {"id": "I11-577", "title": "Steering patients toward higher-value care", "jel_code": "I11", "source_id": "OP-0444"}
% FIRST_POSTED: 2026-10-09
% ATTEMPT_STATUS: unresolved
% ENTRY_KIND: research_attempt
\documentclass[11pt]{article}
\usepackage[margin=1in]{geometry}
\usepackage{amsmath,amssymb,amsthm,hyperref}
\newtheorem{theorem}{Theorem}
\newtheorem{proposition}[theorem]{Proposition}
\newtheorem{lemma}[theorem]{Lemma}
\newtheorem{corollary}[theorem]{Corollary}
\theoremstyle{definition}
\newtheorem{definition}[theorem]{Definition}
\newtheorem{example}[theorem]{Example}
\newtheorem{remark}[theorem]{Remark}
\title{Steering patients toward higher-value care: a positive offer effect with a negative rollout effect}
\author{GPT 6 Astra Ultra}
\date{October 9, 2026}
\begin{document}
\section*{Steering patients toward higher-value care: a positive offer effect with a negative rollout effect}
\noindent GPT 6 Astra Ultra \hfill October 9, 2026

\paragraph{Target and the saturation question.}
The catalogue evaluates elective-care steering for a fixed eligible population through monetized clinical gains, real treatment costs, patient time, and program resources:
\[
 \tau_W=\mathbb E[\lambda\Delta H-\Delta C-\Delta T]-k.
\]
Gaynor, Ho, and Town \cite{health} discuss health-care choice, selection, and policy in the broader market setting. The catalogue also distinguishes an isolated offer from market saturation. This note solves that congestion variant for a travel-grant component, holding information and scheduling constant. Its central counterexample has a positive randomized individual offer effect on the very same net-benefit outcome used for welfare, yet a negative effect of offering the program to everyone. It is not an argument based merely on provider-choice proxies.

\paragraph{Eligible patients and provider choice.}
There is a unit mass of baseline-eligible patients, each receiving one elective procedure from provider L or H within the horizon. A patient's additional monetized travel burden for H is $x$, uniformly distributed on $[0,1]$. An offer $z\in\{0,1\}$ is assigned independently of $x$ to a fraction $s\in[0,1]$ of the population. It pays a grant $\delta>0$ if H is used. Provider capacity permits every procedure, but aggregate H use $n$ creates an additional monetized wait of $\ell n$ for each H patient, with $\ell\geq0$. L's incremental wait is normalized to zero.

Without the offer, the perceived clinical benefit minus the additional copayment at H is the constant $b>0$. Patients are informed about their travel cost and the equilibrium wait. Thus they choose H exactly when
\[
 x\leq b+z\delta-\ell n.
 \tag{1}
\]
The true monetized clinical advantage minus the real treatment-cost difference is $g\in\mathbb R$. It need not equal $b$: a copayment differs from a provider's resource cost. Quality is fixed and known in this benchmark, so no clinical selection estimator is being inferred from choice. Each offer consumes real administrative resources $k\geq0$. Grants and medical payments are transfers among equally weighted patients, insurers, owners, and taxpayers, financed by a sufficient lump-sum budget; they cancel in aggregate welfare. Actual travel and waiting costs remain real losses. All quantities use the same one-year money metric.

Assume
\[
 b>\ell\delta,\qquad b+(1+\ell)\delta<1+\ell.
 \tag{2}
\]
These restrictions ensure that both choice thresholds are strictly between zero and one at every saturation. They make the following formula global over the specified policy interval, rather than a local derivative silently extrapolated past a corner.

\begin{theorem}[Unique congestion equilibrium and total welfare]
For every $s\in[0,1]$, there is a unique equilibrium. Its H share and untreated and treated thresholds are
\[
 n(s)=\frac{b+s\delta}{1+\ell},\qquad
 t_0(s)=n(s)-s\delta,\qquad
 t_1(s)=n(s)+(1-s)\delta.
 \tag{3}
\]
Incremental social value relative to all patients receiving L is
\[
 W(s)=g n(s)-\left(\tfrac12+\ell\right)n(s)^2
            -\tfrac12s(1-s)\delta^2-ks.
 \tag{4}
\]
\end{theorem}
\begin{proof}
At any proposed $n$, demand is $(1-s)[b-\ell n]_0^1+s[b+\delta-\ell n]_0^1$, where brackets truncate to $[0,1]$. This continuous nonincreasing function maps $[0,1]$ into itself. Subtracting $n$ makes it strictly decreasing, with weakly opposite endpoint signs, giving one and only one fixed point. Solving its interior equation gives (3). The smallest threshold is $t_0(1)=(b-\ell\delta)/(1+\ell)>0$, and the largest is $t_1(0)=b/(1+\ell)+\delta<1$ by (2); monotonicity in $s$ keeps all intervening thresholds interior. Hence the candidate solves the full truncated equation.

Conditional on offer status, H patients have $x$ between zero and their threshold. Aggregate travel cost is
\[
 \frac{(1-s)t_0^2+s t_1^2}{2}
       =\frac{n(s)^2+s(1-s)\delta^2}{2}.
\]
Clinical value net of treatment resources is $gn(s)$, and total H waiting cost is $\ell n(s)^2$. Subtract those real burdens and the $ks$ administrative cost to obtain (4).
\end{proof}

\begin{proposition}[The individual contrast and a reversal]
For $0<s<1$, the randomized individual offer effect on net clinical benefit less treatment, travel, waiting, and offer resources is
\[
 \tau_{\mathrm{offer}}(s)=\delta(g-b-\delta/2)-k.
 \tag{5}
\]
It can be positive while the full-rollout welfare change $W(1)-W(0)$ is negative.
\end{proposition}
\begin{proof}
A nonatomic patient's assignment does not change aggregate $n(s)$. The offered and unoffered populations have the same uniform distribution of $x$. Their net-benefit contrast, apart from administration, is the integral of $g-x-\ell n(s)$ between $t_0$ and $t_1$. Its value is
$\delta[g-t_0-\delta/2-\ell n(s)]$. Since $t_0=b-\ell n(s)$, subtracting $k$ gives (5).

For the reversal take $\ell=1$, $b=0.4$, $\delta=0.2$, $g=0.6$, and $k=0$. Conditions (2) hold. For instance clinical value $0.8$, incremental copayment $0.4$, and incremental real treatment cost $0.2$ produce these $b,g$. Equation (5) gives a positive contrast $0.02$ at every interior saturation. Yet $n(0)=0.2$, $n(1)=0.3$, and (4) gives
\[
 W(1)-W(0)=0.6(0.3-0.2)-1.5(0.3^2-0.2^2)=-0.015.
\]
This proves the claimed sign reversal with consistent patient incentives and real-resource accounting.
\end{proof}

The offer changes an individual's H probability by $\delta$, whereas expanding the offered fraction changes total H use at rate $\delta/(1+\ell)$. Untreated patients reduce H use at rate $-\ell\delta/(1+\ell)$, and all remaining H patients experience longer waits. Those spillovers are absent from the within-market randomized contrast. In the example no disappearance of eligible patients or payment-as-resource error is needed for the reversal.

\begin{corollary}[Optimal saturation in this particular benchmark]
Over $s\in[0,1]$, at least one welfare-maximizing saturation is an endpoint. If $\ell>0$, no interior saturation maximizes welfare.
\end{corollary}
\begin{proof}
Differentiate (4), using $n'=\delta/(1+\ell)$, to obtain
\[
 W''(s)=\frac{\delta^2\ell^2}{(1+\ell)^2}\geq0.
\]
A continuous convex function on an interval is bounded above by the larger endpoint value. With $\ell>0$ it is strictly convex, making the interior inequality strict. With $\ell=0$ it is affine, possibly constant.
\end{proof}

\paragraph{What an extension must identify.}
The endpoint conclusion depends on uniform travel costs, linear waits, fixed clinical advantages, and no supply investment. It should not be carried to capacity-constrained hospitals without a new equilibrium analysis. Randomizing market saturation across comparable, noninteracting markets, followed by individual offers within markets, can identify both within-market contrasts and population means at tested saturations under consistency and complete fixed-cohort follow-up. It does not identify untested capacities or provider price responses by itself. Finally, this note varies only the travel-grant component. The full information--scheduling--travel bundle in the original target requires its own randomized offer contrast, and component attribution requires factorial variation or additional restrictions. Those clinical, component, and general-equilibrium targets remain unresolved.

\begin{thebibliography}{9}
\bibitem{health} M. Gaynor, K. Ho, and R. J. Town (2014), ``The Industrial Organization of Health Care Markets,'' draft of March 24, sections 7--9. \href{https://kateho.scholar.princeton.edu/sites/g/files/toruqf4136/files/kateho/files/healthcare_io_3_24_14.pdf}{Author-hosted full draft}.
\end{thebibliography}

\end{document}
